\section{Further questions}\label{sec:questions}

The two constructions give complete extremal statements within their
classes, but do not classify all stable members. For blowups, the
vertex criterion separates this classification problem from the
subspace-rank calculation. For bundles, the subdivision theorem gives
many extremizers without a characterization of every stable twist.

\begin{question}
Which saturated blowup trees of maximum degree four or five have
stable tangent bundle? What additional subspace inequalities are needed
for the same construction on graphs with cycles?
\end{question}

The degree-four obstruction is sufficient, whereas
Proposition~\ref{treemix:two-centres} gives stable higher-degree
arrangements. Larger factor dimensions at nonleaf vertices also lie
beyond the subcubic stability theorem. The exceptional mixed-intersection
formula in Section~\ref{sec:intrinsic} remains available without a tree,
but controlling its associated foliation alone does not control every
higher-rank subsheaf.

\begin{question}
Is there an absolute integer $L_*$ such that every finite subcubic
tree with an edge gives a stable graph bundle whenever all skeleton
edges are subdivided to length at least $L_*$?
\end{question}

Theorem~\ref{tree:stable} gives a sufficient length growing
logarithmically with the number of trivalent vertices. It does not
show that this dependence is necessary. More generally, the necessary
incidence conditions in Corollary~\ref{interact:equality} do not decide
which primitive twisting vectors produce stable extremizers.

Finally, let $R(n)$ be the largest Picard number of a smooth toric
Fano $n$-fold with anticanonically stable tangent bundle. The bundle
theorem gives the lower bound below. For comparison, a published
splitting theorem supplies a general upper bound.

\begin{corollary}[Consequence of Assarf--Nill]\label{cor:publishedbound}
For $n\ge4$,
\[
 \left\lfloor\frac{5n}{4}\right\rfloor+1\le R(n)\le
 2n-\left\lceil\frac{\sqrt{60n+1309}-37}{30}\right\rceil.
\]
\end{corollary}
\begin{proof}
The lower bound is Theorem~\ref{thm:main}. For a stable $X$, put
$k=2n-\rho(X)\ge0$, using Casagrande's bound
\cite{Casagrande2006}. Assarf--Nill
\cite[Theorem~8]{AssarfNill2016} show that
$n>15k^2+37k+1$ forces a factor $S_6$ in a product decomposition.
For $n\ge3$ this is a nontrivial product, whose tangent bundle is
not stable. Thus $n\le15k^2+37k+1$; solving for the nonnegative
integer $k$ gives the upper bound.
\end{proof}

Assarf--Nill \cite[Conjecture~13]{AssarfNill2016} conjecture the
stronger bound $\rho\le(4n+3)/3$ for smooth toric Fano varieties of
dimension at least three that are not nontrivial products. Tangent
stability implies this nonproduct condition.
\begin{question}
Does every smooth toric Fano $n$-fold with anticanonically stable
tangent bundle satisfy $\rho\le(4n+3)/3$ for $n\ge3$? What is
$\limsup_{n\to\infty}R(n)/n$?
\end{question}

The results here give an asymptotic lower ratio of $5/4$. The proposed
upper ratio $4/3$ remains conjectural. A general proof would have to
show that excessive Picard number forces excessive degree in some
proper tangent subspace. Neither the exceptional-intersection identity
nor the restricted bundle incidence argument supplies that implication.
